\chapter{Session procedure}
\label{chap:procedure}

\section{Why the order matters}

The sequence below is not a checklist that can be reordered for convenience.
The geometric settings are fixed before the photometric ones for a specific
reason: \textbf{moving a fixture to correct uniformity changes the
illuminance, while dimming it does not change the uniformity}. Settle
position and aim first, then level. Doing it the other way round means
measuring the light twice.

The same logic places calibration after the optics are locked and before any
experimental capture, and places the stray-light check before everything,
because a room that is not actually dark invalidates every photometric number
measured in it.

\begin{figure}[htbp]
\centering
\begin{tikzpicture}[
  node distance = 6mm and 10mm,
  box/.style  = {draw, rounded corners=2pt, align=center, inner sep=4pt,
                 font=\small, minimum height=8mm, text width=30mm},
  gate/.style = {draw, align=center, inner sep=4pt, font=\small\itshape,
                 minimum height=8mm, text width=30mm},
  ar/.style   = {-{Stealth[length=2mm]}}
]
\node[box] (s1) {1. Darken \\ and check};
\node[box, below=of s1] (s2) {2. Warm up};
\node[box, below=of s2] (s3) {3. Place and aim \\ (geometry)};
\node[gate, below=of s3] (s4) {4. Uniformity \\ grid};
\node[box, below=of s4] (s5) {5. Level and \\ record lux};
\node[box, right=22mm of s1] (s6) {6. Set and lock \\ the optics};
\node[box, below=of s6] (s7) {7. Flicker test};
\node[box, below=of s7] (s8) {8. Calibrate on \\ the support};
\node[box, below=of s8] (s9) {9. Run the \\ experiment};
\node[box, below=of s9] (s10) {10. Check marks, \\ complete record};

\draw[ar] (s1) -- (s2);
\draw[ar] (s2) -- (s3);
\draw[ar] (s3) -- (s4);
\draw[ar] (s4) -- (s5);
\draw[ar] (s5.east) -- ++(6mm,0) |- (s6.west);
\draw[ar] (s6) -- (s7);
\draw[ar] (s7) -- (s8);
\draw[ar] (s8) -- (s9);
\draw[ar] (s9) -- (s10);

\draw[ar] (s4.west) -- ++(-8mm,0) |- node[pos=0.25, left, font=\scriptsize,
          align=right] {fails \\ A2.4} (s3.west);
\end{tikzpicture}
\caption[Session procedure]{The session procedure. The only loop is the
uniformity check returning to fixture placement, because correcting
uniformity means moving a fixture, which changes the illuminance and
invalidates any level already measured.}
\label{fig:procedure}
\end{figure}

\section{The ten steps}

\begin{enumerate}
    \item \textbf{Darken the room.} Close, cover, and rig the black
    background. Verify with a frame from each camera at the longest exposure
    that will be used, with the fixture off. Nothing recognisable may appear.
    Anything visible is stray light and must be eliminated before continuing.
    This is criterion A2.1.

    \item \textbf{Warm up the fixture} for the manufacturer's stated time
    before measuring anything. For HMI this is several minutes. If a
    spectroradiometer is used, EBU recommends at least 30 minutes for the
    instrument itself \cite{ebu_tech3355}.

    \item \textbf{Place and aim.} Position the fixture at the greatest throw
    the room allows and set its aim. Run the penumbra test of Section
    \ref{sec:penumbra} and record the width. This is criterion A2.3.

    \item \textbf{Map the uniformity.} Take a lux grid across width, height
    and depth of the working volume. If criterion A2.4 fails, move the
    fixture and return to step 3. Do not attempt to fix uniformity by
    dimming.

    \item \textbf{Record the level.} Illuminance at the tether position,
    against criterion A2.5. Also record the grid minimum and maximum from
    step 4 and the measured colour temperature.

    \item \textbf{Set and lock the optics.} Aperture wide open, focus set by
    eye on the live full-resolution image with the tether in place, then both
    rings tightened and marked. Photograph the marks.

    \item \textbf{Flicker test.} White card, shortest exposure in use, one
    frame from each camera. Look for banding. This is criterion A2.7. If it
    fails, the fixture or its ballast mode has to change, and steps 3 to 5
    are repeated for the replacement.

    \item \textbf{Calibrate} on the integrated support, in place, following
    Section \ref{sec:calib-run}.

    \item \textbf{Run the experiment} following Section \ref{sec:capture-plan},
    touching nothing on the optics.

    \item \textbf{Close out.} Check the ring marks against the photographs
    from step 6 and record whether they moved. Complete every field of the
    record sheet. Copy the captures off the device.
\end{enumerate}

\section{Running the calibration}
\label{sec:calib-run}

The application's guided flow does the bookkeeping. The operator's job is
physical.

The tool asks for one pose at a time and draws its outline in both camera
previews. A frame is captured only when it is sharp over the board region,
momentarily still, visible to \emph{both} cameras, and on the pose currently
being asked for. A rejected frame does not consume a slot, so a request for
25 views yields 25 that passed every gate.

Three practical points for the operator:

\begin{itemize}
    \item \textbf{Stand back and hold the board yourself.} The flow is built
    for a single person holding the board at a distance and matching the
    drawn outline. The preview is sized so it remains readable from across
    the room.

    \item \textbf{The strongly tilted poses are the point.} The bins at 30 to
    45\textdegree{} and 45 to 70\textdegree{} are the ones no archived session
    ever reached, and they are the ones that make the focal length
    recoverable. If a pose cannot be physically reached, use \emph{Finish with
    the views collected}; that is the deliberate way out, and the session
    reports itself as incomplete rather than pretending otherwise.

    \item \textbf{Watch the uncertainty, not the residual.} The progress
    indicator is driven by parameter uncertainty from the fit covariance, not
    by a frame count and not by the reprojection error. A session is finished
    when the uncertainty stops falling, not when a number of frames has been
    reached.
\end{itemize}

At the end, evaluate Table \ref{tab:crit-calib}. If A1 fails, the experiment
can still be run, but Chapter \ref{chap:analysis} will report lengths in
pixels and ratios rather than in metres, and criterion A3.5 is void.

\section{The capture plan}
\label{sec:capture-plan}

Both the tether and the rig are moved by hand. The plan is organised so that
the simplest configurations are captured first, while everything is known to
be working, and the hardest last.

\begin{table}[htbp]
\centering
\caption[Capture plan]{Capture sequence. Each configuration is captured as
several stereo pairs, so that a failure in any one pair does not cost the
configuration. Counts are a minimum.}
\label{tab:capture-plan}
\begin{tabularx}{\textwidth}{@{}l X l@{}}
\toprule
\textbf{Set} & \textbf{Configuration} & \textbf{Pairs} \\
\midrule
S1 & Tether taut and static, whole length in frame in both cameras. The
     reference set for A3.5 and A3.6 & $\geq 5$ \\
S2 & Tether with deliberate sag, static & $\geq 3$ \\
S3 & Tether static, camera rig moved slowly by hand & $\geq 5$ \\
S4 & Tether static, camera rig moved erratically & $\geq 5$ \\
S5 & Tether in motion, rig static & $\geq 5$ \\
S6 & Both in motion, erratic. The nominal case & $\geq 10$ \\
S7 & Tether arranged with a deliberate self-crossing, static & $\geq 3$ \\
S8 & One endpoint deliberately out of frame in one camera only &
     $\geq 3$ \\
\bottomrule
\end{tabularx}
\end{table}

\textbf{S1 is the measurement set.} It is the only one where an external
reference exists, because the tether is taut and its length has been
tape-measured. Everything else is an exercise of robustness.

\textbf{S7 exists to characterise the tracker, not to pass.} A self-crossing
is where the automatic tracker has to choose a branch using a hand-tuned
scoring formula, and it is the case the project expects to be weakest.
Capturing it deliberately turns an anticipated weakness into a measurement,
and the pairs feed the path dataset whether the tracker succeeds or not.

\textbf{S8 exists to provoke a silent failure.} The parametric matching
assumption breaks when the two traced paths do not cover the same physical
extent, and nothing in the pipeline detects it. Deliberately occluding one
endpoint in one camera creates that condition on purpose, so that the size
and the sign of the resulting error can be seen for the first time against
the tape-measured length from S1.
